% S3: Practice Survey.
%
% Four things are bound here because none of them is read sequentially
% against the chapters that summarise or depend on them: the RQ5 survey
% instrument's specification, including its inclusion and exclusion rule
% and its strict/permissive coding convention; the survey's results, the
% coded-outcome table and the reliability statistics behind it; the
% limitation the discussion chapter states about this measurement; and the
% full 75-paper coding table itself, bound in full immediately afterwards so
% a reader can recompute the counts rather than take them on trust. Numbers
% and claims are unchanged from the thesis chapters and the appendix file
% they were moved from; only cross-references have been adjusted for the
% move, as noted at each site.

\section{Practice Survey}
\label{sup:s3}

\subsection{RQ5: the survey coding instrument}

\textbf{Condition.} A survey of practice is evidence about practice only if the coding rule is
fixed, applied uniformly, and reported under both the strict and the permissive reading of every
contestable boundary. Reporting only the reading that supports the thesis is the failure the survey
exists to expose.

The instrument distinguishes two audits a paper may perform. A model audit examines a trained model
for disparities across groups within the analysis pool. A selection audit examines the filtering
step itself for differential retention across groups, against a denominator existing prior to
filtering. Inclusion requires all three of: health, clinical or biomedical data about people; at
least one disaggregated performance or outcome metric broken out by a demographic or socioeconomic
attribute; and enough methods detail to determine what cohort the metric was computed on. Papers
are excluded when they are commentary, position or review articles with no empirical evaluation,
when they report no disaggregated metric, when the data is not health data, when the methods are
too thin to identify the cohort, or when the full text is not retrievable.

Queries were issued against machine-readable interfaces (NCBI E-utilities for PubMed and PubMed
Central, the arXiv Atom API, and ACL Anthology volume indexes), so the search is reproducible by
re-running the recorded query strings. Screening runs in three stages: an automated title and
abstract screen logging its exact exclusion rule for each record; a full-text retrieval quota
ranking records passing stage one by a pre-specified keyword-density score and advancing them until
the budget is met, with records below the quota logged as screened rather than silently dropped;
and a full-text screen read against the three inclusion criteria. The corpus reaching coding
contains $n = 75$ papers.

The primary coding of a selection audit is strict: a paper counts only if it compares the group
composition, or the retention rate, of the retained cohort against a denominator existing prior to
filtering, or otherwise directly tests the filter for differential retention; a demographics table
or a flow diagram of total counts excluded does not qualify, since neither has a denominator. A
second, permissive column counts any comparison of retained group composition against any external
or source figure, including census statistics or the parent database before filtering; it is a
strict superset of the strict column, enforced in code, and both readings are reported everywhere.
Every coded field requires a verbatim quote and a section location, checked by a verification
script against the source text by exact string match, and a quote that fails to match is rejected
and its field forced to unclear.

Three limitations bound what the instrument can measure. Each paper was coded once by the primary
coder under a shared written rule set; a blind second coder independently re-coded a deterministic
subsample of nineteen papers, the basis for the agreement statistics reported in
Section~\ref{sup:sec:survey-reliability}, and no agreement figure is claimed for the other fifty-six. Second,
inclusion criterion two requires a disaggregated metric, so every included paper audits a model on
its pool by construction: the count of papers coded as auditing a model on the pool equals the count
of included papers, both $75$, so the cross-tabulation of model audit against selection audit
degenerates to a single row, and the keyword-density ranking enriches the corpus toward
fairness-forward work, against the direction of the finding rather than toward it. Third, the coded
question is only whether a paper audits the selection step upstream of the model it evaluates. It is
not the exemplar-correctness check, the weak-label marginal check, the architecture randomisation
control, or the positive control for evaluation sensitivity, and the survey measures none of them,
so the corpus supports only the prevalence of this one adjacent practice, not a claim about whether
the field runs the controls that would have caught this pipeline's own failures.



\subsection{Practice: RQ5 Survey Results}
\label{sup:sec:res-rq5}

Does the field run the checks that would have caught the failures above? This
section reports a prevalence measurement on a deliberately enriched corpus, not
a finding of the same standing as the ablation in
S1 of this supplement or the pre-registered positive control in
Supplement S6. It is reported because a prevalence measurement on
its own terms is worth having, and because the survey coding table below makes
the full 75-paper coding available for a reader to check rather than take on
trust. Of the $525$ coded fields, $502$ carry a verbatim quote checked by exact string match
against the source text, so no coded value asserting a property of a paper survives without
verified textual support. The remaining $23$ are codings of absence, recording that a paper does
not mention how it handles null attribute values; there is no sentence to quote for an absence,
and those $23$ are the \emph{not mentioned} row of Table~\ref{sup:tab:survey}. A systematic coding of 75 papers says: rarely on the strict
reading, and not much more often on the permissive one. Both are reported side by side; the
choice between them is a judgement call for the reader.

\subsection{Both readings, and the null-attribute distribution}

\begin{table}[H]
\caption[Coded outcomes across surveyed papers]{Coded outcomes across 75 included papers.}
\centering
\small
\begin{tabular}{lr}
\toprule
\textbf{Coded outcome} & \textbf{Papers} \\
\midrule
Runs a model over the candidate pool  & 75 \\
Reports the check, permissive reading & 14 \\
Reports the check, strict reading     & 4  \\
\midrule
\multicolumn{2}{l}{\textit{Handling of null attribute values}} \\
Not mentioned     & 23 \\
Drops the rows    & 17 \\
Separate category & 13 \\
Unclear           & 9  \\
Imputes           & 7  \\
Reports only      & 6  \\
\bottomrule
\end{tabular}
\apanote{Papers are coded once each; the denominator is 75 throughout. The first row is the screening condition, satisfied by every included paper by construction. The permissive reading accepts any statement that the check was performed; the strict reading requires the quantity that would make it verifiable. The six null-attribute categories are exclusive and sum to 75.}
\label{sup:tab:survey}
\end{table}

Four of 75 papers report the check under the strict reading and 14 under the
permissive one, so 61 report it under neither. The 10-paper gap between the two
readings is the population that says it did the check without reporting what the
check returned, and the strict count is the one a reader can verify. Twenty-three
papers, the largest single category, do not mention the null-attribute question
and nine more are coded unclear, so 32 of 75 give a reader no way to tell what
happened to rows with a missing attribute. Dropping and imputing are not
equivalent and neither is neutral, and the two choices move a reported disparity
in different directions; 17 papers take the first and 7 the second. The point is
not that one choice is wrong; it is that 32 of 75 leave it unrecoverable.

\subsection{Reliability}
\label{sup:sec:survey-reliability}

One coder coded all 75 papers under the shared written rule set described in
the coding-instrument subsection above. A blind second coder independently re-coded a
19-paper subsample, given the field definitions but not the hypothesis, and
that subsample is the basis for every agreement figure in this section; no
agreement is claimed for the remaining $75 - 19 = 56$ papers, which were coded
once. On the permissive cell agreement is $0.9474$ and Cohen's kappa is
$0.855$; on the
null-attribute-handling coding agreement is $0.7895$ and kappa is $0.730$.

The strict cell has no defined agreement coefficient, for a structural
reason: chance-corrected agreement
\citep{cohen1960kappa} is a ratio whose denominator is one minus expected
agreement, and when a coder's marginal on a category is degenerate that
denominator goes to zero and the coefficient is undefined rather than perfect. At
4 positive codes in 75 the strict cell is degenerate on the double-coded
subsample specifically: on the 19 overlap papers neither coder found a
strict-positive paper, so agreement is $1.0$ but kappa is undefined rather
than perfect. Reporting perfect agreement as a defined kappa would be exactly
the error the statistic exists to prevent.

The corpus is enriched, and the enrichment runs against the finding rather
than for it: a corpus that over-samples work already attentive to these
questions still yields 4 of 75 strictly, which bounds the bias without
turning 75 papers into a population estimate for the field.


\paragraph{The survey is a prevalence measurement on an enriched corpus, and it
is limited on several axes at once.} It is the weakest of the empirical
strands in this thesis, reported at that weight rather than alongside the
ablation or the pre-registered convergence control. The $75$-paper corpus was
assembled to over-sample fairness-forward work, coverage is open-access
sources only, $459$ screened-in papers were never retrieved in full text, and
one coder codes the full set; the reliability subsection above gives the
enrichment argument behind $4/75$ and $14/75$. A blind second coder recoded a
subsample with field definitions only and the hypothesis withheld, and
agreement on the permissive field is high. On the 19 double-coded papers, the
registered permissive
agreement is $0.9474$ with kappa $0.855$. On the strict cell, the one that decides the
headline count, the reliability coefficient is
\textbf{undefined}: both coders returned the same code on all 19 overlap
papers, which drives a coder's marginal on the category toward degenerate and
the denominator of a chance-corrected coefficient toward zero. The full
contingency table behind the undefinedness is not registered and is not
claimed here beyond that. It is reported as undefined rather than as the
perfect agreement the raw agreement rate superficially permits, which is the
same class of error as reporting a coefficient guaranteed by construction.
The survey coding table below reproduces the full 75-paper coding, so a
reader who wants to check the count rather than take it on trust can.


\subsection{Survey Coding Table}

The full coding behind the counts above, reproduced from the appendix that
has carried it since it was first bound.

\begingroup
% msu_appD_survey.tex opens with \chapter{...}\label{app:survey}. This
% document has no chapter-lettered appendices to house it in, so \chapter is
% suppressed for the extent of this input; the subsection heading above
% already gives the reader the heading \chapter would otherwise have
% printed. \label{app:survey} still fires, against whatever section-level
% \@currentlabel is active at that point, and is not relied on by any \ref
% in this document.
\renewcommand{\chapter}[1]{}
\chapter{Survey Coding Table}
\label{app:survey}

The prevalence measurement in the results chapter of the thesis rests on 75 papers coded
against a written instrument. The full coding is reproduced here so the counts can be checked
rather than taken on trust. Under the strict reading
4 of 75 papers audit the selection step; under the permissive reading
14 do. Every field value in the underlying record carries a verbatim quote with its section
location, 225 in total, and those quotes were verified as exact substrings of the retrieved
text. The quotes themselves, the screening log including excluded and unretrievable papers, and
the coding protocol are released alongside this thesis.

\singlespacing
\begin{longtable}{@{}p{0.9in}p{0.42in}p{0.42in}p{0.55in}p{1.7in}@{}}
\textbf{Paper} & \textbf{Model audit} & \textbf{Filtered cohort} & \textbf{Audits selection} & \textbf{Null attribute handling}\\
\midrule
\endhead
AC2020\_clinicalnlp-1\_33 & yes & yes & no & separate\ category\\
AC2022\_clinicalnlp-1\_10 & yes & yes & no & not\ mentioned\\
AC2023\_bionlp-1\_39 & yes & yes & no & unclear\\
AC2026\_bionlp-1\_14 & yes & yes & no & not\ mentioned\\
AX2003.00827 & yes & yes & no & drops\\
AX2003.11515 & yes & yes & no & drops\\
AX2111.09507 & yes & yes & no & reports\ only\\
AX2304.02780 & yes & yes & no & unclear\\
AX2304.05986 & yes & yes & no & not\ mentioned\\
AX2308.02588 & yes & yes & no (yes) & drops\\
AX2308.06354 & yes & yes & no & not\ mentioned\\
AX2312.06979 & yes & yes & no & not\ mentioned\\
AX2410.01089 & yes & yes & no & not\ mentioned\\
AX2411.08048 & yes & yes & no & separate\ category\\
AX2412.20374 & yes & yes & no & not\ mentioned\\
AX2502.09626 & yes & yes & no & not\ mentioned\\
AX2502.16051 & yes & yes & no & not\ mentioned\\
AX2503.15454 & yes & yes & no & not\ mentioned\\
AX2509.09772 & yes & yes & no & imputes\\
AX2509.15107 & yes & yes & no & drops\\
AX2510.00055 & yes & yes & no (yes) & not\ mentioned\\
AX2511.15986 & yes & yes & no & separate\ category\\
AX2512.19735 & yes & yes & no & unclear\\
AX2602.20303 & yes & yes & no & imputes\\
AX2603.14631 & yes & yes & no (yes) & drops\\
AX2604.04858 & yes & yes & no & unclear\\
AX2605.04882 & yes & yes & no & not\ mentioned\\
AX2605.19214 & yes & yes & no & not\ mentioned\\
AX2607.16253 & yes & yes & no (yes) & reports\ only\\
AX2607.29082 & yes & yes & no & drops\\
PM32064914 & yes & yes & yes & drops\\
PM33427747 & yes & yes & no & not\ mentioned\\
PM34798329 & yes & yes & no & drops\\
PM35033309 & yes & yes & no & imputes\\
PM35504931 & yes & yes & no & drops\\
PM35710407 & yes & yes & yes & separate\ category\\
PM36714611 & yes & yes & no & unclear\\
PM36716365 & yes & yes & no & separate\ category\\
PM37318804 & yes & yes & no & reports\ only\\
PM37921762 & yes & yes & no (yes) & unclear\\
PM38422347 & yes & yes & no & not\ mentioned\\
PM38985468 & yes & yes & no & separate\ category\\
PM39153563 & yes & yes & no & drops\\
PM39386055 & yes & yes & no & reports\ only\\
PM39399154 & yes & yes & no (yes) & unclear\\
PM39625723 & yes & yes & no & separate\ category\\
PM39671594 & yes & yes & no (yes) & not\ mentioned\\
PM39825353 & yes & yes & no & drops\\
PM39831110 & yes & yes & no & separate\ category\\
PM39942473 & yes & yes & no & not\ mentioned\\
PM40047796 & yes & yes & no & drops\\
PM40111994 & yes & yes & no & drops\\
PM40298901 & yes & yes & no & separate\ category\\
PM40598240 & yes & yes & no & reports\ only\\
PM40639839 & yes & yes & no & drops\\
PM40799932 & yes & yes & no & unclear\\
PM40824638 & yes & yes & no & separate\ category\\
PM40841556 & yes & yes & no & not\ mentioned\\
PM40866555 & yes & yes & no & not\ mentioned\\
PM40893625 & yes & yes & no & drops\\
PM40905712 & yes & yes & no & not\ mentioned\\
PM40971151 & yes & yes & yes & separate\ category\\
PM41351059 & yes & yes & no & imputes\\
PM41487722 & yes & yes & no & unclear\\
PM41639760 & yes & yes & no & not\ mentioned\\
PM41653838 & yes & yes & no & drops\\
PM41674602 & yes & yes & no (yes) & not\ mentioned\\
PM41772035 & yes & yes & no & drops\\
PM41840801 & yes & yes & no & imputes\\
PM42078370 & yes & yes & no & separate\ category\\
PM42081529 & yes & yes & no (yes) & imputes\\
PM42210248 & yes & yes & yes & separate\ category\\
PM42243352 & yes & yes & no (yes) & imputes\\
PM42434194 & yes & yes & no & reports\ only\\
PM42500108 & yes & yes & no & not\ mentioned\\
\end{longtable}

\apanote{Values in the audits-selection column are the strict reading; a parenthesised value is the permissive reading where the two differ, which is the ten-paper band the thesis reports as the honest ambiguity in this measurement.}
\doublespacing


\endgroup
